\documentclass[journal, 10pt]{IEEEtran}
\usepackage{amsmath, amsfonts, amssymb}
\usepackage{graphicx}
\usepackage{bm}
\usepackage{cite}
\usepackage{hyperref}

\title{Solid-State Steerable Synthetic Jet Actuator via Phased-Array Acoustic Streaming}
\author{Mohana Rangan Desigan\\ \textit{Mac Mini M4 Compute Node, Independent Research}}

\begin{document}
\maketitle

\begin{abstract}
This paper establishes the theoretical baseline and digital twin methodology for a solid-state synthetic jet actuator. Utilizing a 64-element uniform rectangular array (URA) of 40\,kHz piezoelectric transducers, macroscopic fluid flow is generated entirely through the nonlinear attenuation of finite-amplitude acoustic waves. By coupling the Lighthill-Nyborg acoustic streaming equations with phased-array beamforming delay laws, we define a computational pipeline that extracts the time-averaged Reynolds stress tensor from a linear Helmholtz acoustic solution and applies it as a steady-state momentum source in the incompressible Navier-Stokes equations.
\end{abstract}

\section{Introduction \& System Architecture}
Traditional fluid boundary-layer control systems, such as Dielectric Barrier Discharge (DBD) plasma actuators, suffer from high-voltage requirements, ozone generation, and environmental instability. We propose an alternative exploiting Eckart bulk acoustic streaming. 

The physical hardware is constrained to an $8 \times 8$ planar array of 40\,kHz transducers with a pitch of $d = 10.5\,\text{mm}$. Driven at $V_{DD} = 24\,\text{V}$, the system acts as a solid-state synthetic jet. Because the array pitch exceeds the spatial Nyquist criterion ($\lambda/2 \approx 4.28\,\text{mm}$), grating lobes are unavoidable, necessitating strict spatial targeting. Empirical baseline measurements suggest maximum achievable time-averaged flow velocities on the order of $0.1\,\text{m/s}$ to $0.5\,\text{m/s}$ at high Sound Pressure Levels (SPL).

\section{Phased-Array Focusing Mathematics}
To synthesize a focused momentum jet at an arbitrary spatial coordinate $\mathbf{r}_f = (x_f, y_f, z_f)$, the acoustic wavefronts from all $N$ elements must arrive perfectly in phase. Let the position of the $i$-th transducer be $\mathbf{r}_i = (x_i, y_i, 0)$.

The geometric distance from element $i$ to the focal point is:
\begin{equation}
    R_i = \|\mathbf{r}_i - \mathbf{r}_f\| = \sqrt{(x_i - x_f)^2 + (y_i - y_f)^2 + z_f^2}
\end{equation}
To ensure simultaneous wavefront arrival, we define the time-delay law $\Delta t_i$ relative to the element with the shortest path:
\begin{equation}
    \Delta t_i = \frac{R_i}{c_0} - \frac{\min_j R_j}{c_0}
\end{equation}
where $c_0 \approx 343\,\text{m/s}$ is the local speed of sound in air.

Due to spatial aliasing ($d > \lambda/2$), parasitic intensity maxima (grating lobes) will appear at angles $\theta_k$ satisfying:
\begin{equation}
    \sin\theta_k = \frac{k\lambda}{d} - \sin\theta, \quad k=\pm 1, \pm 2, \dots
\end{equation}
These lobes will consume a fraction of the acoustic energy, implicitly bounding the maximum magnitude of the resulting streaming velocity.

\section{Acoustic Streaming Theory}
The generation of macroscopic wind from acoustic waves is a second-order nonlinear phenomenon. We decouple the physics using perturbation theory, where the instantaneous velocity $\mathbf{u}$ and density $\rho$ are expanded as:
\begin{align}
    \mathbf{u} &= \epsilon \mathbf{u}_1 + \epsilon^2 \mathbf{u}_2 \\
    \rho &= \rho_0 + \epsilon \rho_1 + \epsilon^2 \rho_2
\end{align}
where subscript $1$ denotes the first-order (40\,kHz) acoustic field, and subscript $2$ denotes the steady (DC) fluid flow.

As derived by Nyborg and Lighthill, time-averaging the compressible Navier-Stokes equations over one acoustic period reveals that the first-order momentum flux leaves a non-zero residual force. The time-averaged second-order momentum equation behaves as a Stokes flow driven by a volumetric body force density $\mathbf{F}_{\text{rad}}$:
\begin{equation}
    \mu\nabla^2\langle \mathbf{u}_2 \rangle - \nabla \langle p_2 \rangle + \mathbf{F}_{\text{rad}} = 0
\end{equation}
The acoustic radiation force $\mathbf{F}_{\text{rad}}$ is defined by the spatial divergence of the Reynolds stress tensor of the first-order velocity field:
\begin{equation}
    F_{\text{rad},i} = -\frac{\partial}{\partial x_j} \bigl\langle \rho_0 u_{1,i} u_{1,j} \bigr\rangle
\end{equation}

\section{Conclusion \& Next Steps}
With the theoretical framework defined, the next phase bridges the analytical mathematics to finite-volume simulations. The time-delay equations (Eq. 2) will inform the boundary conditions of a linear Helmholtz solver to yield $\mathbf{u}_1$. Subsequently, Eq. 6 will be computed discretely to generate $\mathbf{F}_{\text{rad}}$, serving as the momentum source term in an incompressible Navier-Stokes solver.

\section{Computational Fluid Dynamics \& Navier-Stokes Integration}
\subsection{Numerical Discretization and Dynamic Compilation}
To translate the derived acoustic radiation force $\mathbf{F}_{\text{rad}}$ into a macroscopic kinematic field, the incompressible, transient Navier-Stokes (NS) equations were solved using OpenFOAM (\textit{pimpleFoam}). The governing momentum equation is defined as:
\begin{equation}
    \frac{\partial \mathbf{U}}{\partial t} + \nabla \cdot (\mathbf{U} \otimes \mathbf{U}) - \nabla \cdot (\nu_{\text{eff}} \nabla \mathbf{U}) = -\nabla p + \frac{\mathbf{F}_{\text{rad}}}{\rho_0}
\end{equation}
where $\nu_{\text{eff}}$ is the kinematic viscosity and $\mathbf{U}$ is the macroscopic velocity vector. 

A high-fidelity structured hexahedral mesh of 512,000 cells was generated. To bypass the limitations of standard GUI-based solvers, a custom C++ memory hook was written using OpenFOAM's \texttt{vectorCodedSource}. This allowed the decoupled Reynolds stress tensor, calculated via Python from the ElmerFEM Helmholtz solution, to be injected directly into the NS momentum matrix at runtime. The solver utilized adaptive time-stepping, dynamically adjusting $\Delta t$ to maintain a Courant-Friedrichs-Lewy (CFL) number of $\text{Co} \le 0.95$. 

Running natively on an Apple M4 architecture with 16GB unified memory, the solver achieved matrix convergence (residuals $<10^{-6}$) in approximately 15 minutes of compute time for 100 high-resolution frames.

\subsection{Physical Limits and Shockwave Bounding}
The computational model was initially subjected to an amplified force scaling to evaluate the mathematical limits of the algorithm. By scaling the virtual acoustic pressure, the simulation yielded a highly concentrated, steady-state synthetic jet velocity of $6.98\,\text{m/s}$. 

While mathematically stable, evaluating this result against real-world electro-acoustic limitations provides crucial engineering boundaries. The acoustic streaming force scales with the square of the acoustic pressure ($F \propto p^2$). To generate a $6.98\,\text{m/s}$ jet, the theoretical sound pressure level (SPL) at the focal point must approach $181\,\text{dB}$ ($\approx 22,300\,\text{Pa}$). 

In atmospheric conditions, such extreme pressures exit the linear acoustic regime and induce shockwave steepening due to nonlinear wave propagation. Furthermore, driving a piezoelectric transducer to synthesize $181\,\text{dB}$ would exceed the mechanical fracture limits of standard PZT ceramics. Therefore, we constrain the physical operational target of the hardware to a localized SPL of $160\,\text{dB}$, which scales the expected empirical jet velocity to a physically sustainable $0.5\,\text{m/s}$ to $1.5\,\text{m/s}$.

\section{Electrodynamic (E\&M) Simulation Requirements}
The mechanical success of the synthetic jet is entirely dependent on the structural integrity of the electrical power network. A piezoelectric array of 64 transducers acts as a massive bulk capacitor ($C_{\text{total}} \approx 153\,\text{nF}$). 

The instantaneous current $I$ required to charge this capacitive array to $24\,\text{V}$ at a switching interval $dt \approx 25\,\text{ns}$ is governed by:
\begin{equation}
    I(t) = C_{\text{total}} \frac{dV}{dt} + \frac{V}{R}
\end{equation}
which demands transient current spikes exceeding $140\,\text{A}$. If unmitigated, the equivalent series inductance (ESL) of the PCB traces will trigger catastrophic $Ldi/dt$ voltage overshoots, leading to logic brown-outs and thermal destruction of the MOSFET gate drivers.

Before executing the PCB layout in KiCad, a rigorous Electrodynamic (SPICE) Digital Twin must be constructed. This E\&M simulation will mathematically derive the critical damping resistance ($R_s$), the precise bulk electrolytic capacitance required to prevent 24V supply droop, and the high-frequency ceramic decoupling matrix needed to supply the transient switching current.
\end{document}